%Dokumenteinstellungen und Anpassungen
%Dokumentenklasse "scrbook" - Erweitert um den Verweis auf die Verzeichnisse und Texteigenschaften
\documentclass[12pt, a4paper, oneside, parskip=half, listof=totoc, bibliography=totoc, numbers=noendperiod,ngerman]{scrartcl}


%Anpassung der Seitenränder (Standard bottom ca. 52mm anbzüglich von ca. 4mm für die nach oben rechts gewanderte Seitenzahl)
\usepackage[bottom=48mm,left=25mm,right=25mm]{geometry}

%Tweaks für scrbook
\usepackage{scrhack}

%Blindtext
\usepackage{blindtext}

%Erlaubt unteranderem Umbrücke captions
\usepackage{caption}

%Stichwortverzeichnis
\usepackage{imakeidx}
% Frame figure
\usepackage[export]{adjustbox}
%Kompakte Listen
\usepackage{paralist}

%Zitate besser formatieren und darstellen
\usepackage{epigraph}

%Tabeln
\usepackage{booktabs}

\usepackage{amsmath}
%For Comments duh
\usepackage{comment}

%Glossar, Stichworverzeichnis (Akronyme werden als eigene Liste aufgeführt)
\usepackage[toc, acronym]{glossaries} 
%\usepackage{physics}
%Anpassung von Kopf- und Fußzeile
%beinflusst die erste Seite des Kapitels
\usepackage[automark,headsepline]{scrlayer-scrpage}
\automark{section}
\ihead{\leftmark}
\chead{}
\ohead{\thepage}
\ifoot*{}
\cfoot[\thepage]{}
\cfoot*{}
\ofoot*{\thepage}
\ofoot{}
\pagestyle{scrheadings}


%Auskommentieren für die Verkleinerung des vertikalen Abstandes eines neuen Kapitels
%\renewcommand*{\chapterheadstartvskip}{\vspace*{.25\baselineskip}}

%Zeilenabstand 1,5
\usepackage[onehalfspacing]{setspace}

%Verbesserte Darstellung der Buchstaben zueinander
\usepackage[stretch=10]{microtype}

%Deutsche Bezeichnungen für angezeigte Namen (z.B. Innhaltsverzeichnis etc.)
\usepackage[ngerman]{babel}

%Unterstützung von Umlauten und anderen Sonderzeichen (UTF-8)
\usepackage{lmodern}
\usepackage[utf8]{luainputenc}
\usepackage[T1]{fontenc}

%Einfachere Zitate
\usepackage{epigraph}

%Unterstützung der H positionierung (keine automatische Verschiebung eingefügter Elemente)
\usepackage{float} 

%Erlaubt Umbrüche innerhalb von Tabellen
\usepackage{tabularx}

%Erlaubt Seitenumbrüche innerhalb von Tabellen
\usepackage{longtable}

%Erlaubt die Darstellung von Sourcecode mit Highlighting
\usepackage{listings}

% Für vektor unter zeichen
\usepackage{accents}

%Definierung eigener Farben bei nutzung eines selbst vergebene Namens
\usepackage[table,xcdraw]{xcolor}

%Tikz für Statemaschine drawing
\usepackage{tikz}
\usetikzlibrary{automata, positioning, arrows,arrows.meta,shapes.geometric}
\tikzset{
	->, %makes the edges directed
	>=stealth', % makes the arrow heads bold
	node distance=3cm, % specifies the minimum distance between two nodes. Change if necessary.
	every state/.style={thick, fill=gray!10}, % sets the properties for each ’state’ node
	initial text=$ $, % sets the text that appears on the start arrow	
}

%Grafiken (wie jpg, png, etc.)
\usepackage{graphicx}
%Figure wraping duh
\usepackage{wrapfig}
\usepackage{makecell}
\usepackage{siunitx}
\usepackage{mathtools}
\usepackage{subcaption}
%Grafiken von Text umlaufen lassen
\usepackage{wrapfig}

%Ermöglicht Verknüpfungen innerhalb des Dokumentes (e.g. for PDF), Links werden durch "hidelink" nicht explizit hervorgehoben
\usepackage[hidelinks,german]{hyperref}
\usepackage{cleveref}


%Einbindung und Verwaltung von Literaturverzeichnissen
\usepackage{csquotes} %wird von biber benötigt
\usepackage[style=ieee, backend=biber, bibencoding=ascii]{biblatex}
\addbibresource{references/references.bib}

%Anpassung der Überschriften
\addtokomafont{disposition}{\rmfamily}

%Zusätzliche Farben
\definecolor{darkgreen}{RGB}{0,100,0}

%Umbenennungen
\renewcommand{\lstlistlistingname}{Quelltextverzeichnis}

%Pluszeichen in der Referenc beim zitieren ausblenden
\renewcommand*{\labelalphaothers}{}

%Anpassugen zur Quelltextdarstellung, kann bei Bedarf überschrieben werden (z.B. wenn unterschiedliche Sprachen zum Einsatz kommen)
\renewcommand{\lstlistingname}{Codeauszug}
\lstset{
	language=Java,
	numbers=left,
	columns=fullflexible,
	aboveskip=5pt,
	belowskip=10pt,
	basicstyle=\small\ttfamily,
	backgroundcolor=\color{black!5},
	commentstyle=\color{darkgreen},
	keywordstyle=\color{blue},
	stringstyle=\color{gray},
	showspaces=false,
	showstringspaces=false,
	showtabs=false,
	xleftmargin=16pt,
	xrightmargin=0pt,
	framesep=5pt,
	framerule=3pt,
	frame=leftline,
	rulecolor=\color{green},
	tabsize=2,
	breaklines=true,
	breakatwhitespace=true,
	prebreak={\mbox{$\hookleftarrow$}}
}

%Anpassungen für das Abkürzungsverzeichnis
\newglossarystyle{dottedlocations}{%
	\renewcommand*{\glossaryentryfield}[5]{%
		\item[\glsentryitem{##1}\glstarget{##1}{##2}] \emph{##3}%
		\unskip\leaders\hbox to 2.9mm{\hss.}\hfill##5}%
	\renewcommand*{\glsgroupskip}{}%
}

%Titelformen - gewünschtes Layout einkommentieren

%%Graduation
\makeatletter

\newcommand*{\gradeType}[1]{\gdef\@gradeType{#1}}
\newcommand*{\firstExaminer}[1]{\gdef\@firstExaminer{#1}}
\newcommand*{\secondExaminer}[1]{\gdef\@secondExaminer{#1}}
\newcommand*{\matrikelnr}[1]{\gdef\@matrikelnr{#1}}
\newcommand*{\submitDate}[1]{\gdef\@submitDate{#1}}


\renewcommand*{\maketitle}{
	\begin{titlepage}
				\newgeometry{left=2.5cm,right=2.5cm,top=3.0cm,bottom=2.5cm}
		
		\begin{minipage}{0.5\textwidth}
			\includegraphics[height=0.1\textheight]{./Bilder/FB1_ET_rgb}
		\end{minipage}
		\begin{minipage}{0.5\textwidth}
			\begin{flushright}
				\includegraphics[height=0.1\textheight]{./Bilder/HTW_Logo_rgb}
			\end{flushright}
		\end{minipage}
		\begin{center}
			\vfill
			{\Large \@title\par}
			\vskip 0.5cm
			{\large \bfseries Praktikumsbericht\par}
			\vskip 0.5cm
			{\large zur Erlangung der Prüfungsvorleistung im Fach\\ \bfseries \@gradeType}
			\vskip 0.5cm
			{\large an der}
			\vskip 0.5cm
			{\large Hochschule für Technik und Wirtschaft Berlin\\ Fachbereich I - Energie und Information \\ Studiengang Elektrotechnik}
			\vfill
			\begin{flushleft}
				\begin{tabular}[t]{rl}
					1. Prüfer: &\@firstExaminer\\
					%2. Prüfer: & \@secondExaminer\\
					\\
					Eingereicht von: &\@author\\
					Matrikelnummer: & \@matrikelnr\\
					Datum der Abgabe: & \@submitDate
				\end{tabular}
			\end{flushleft}
		\end{center}
		\restoregeometry
	\end{titlepage}
}
\makeatother
\gradeType{Vertiefung Leistungselektronik}
\secondExaminer{Max Mustermann}

%Research paper
%\include{titles/research_papger}
%\subTitle{Ein optionaler Untertitel der Arbeit}
%\researchPart{A}

%Angaben zur Arbeit und dem Author (von beiden Layouts genutzt)
\title{Laborbericht:  Undeland-Beschaltung}
\author{Milan Daniel Larsen}
\matrikelnr{s0581929}
\submitDate{19.10.2022}
\firstExaminer{Prof. Dr.-Ing. J. Ranneberg}

%Verzeichnisse generieren
\makeglossaries
\loadglsentries{references/glossary_acronyms.tex}
\setacronymstyle{long-short}

\makeindex[columns=2, title=Stichwortverzeichnis, options= -s resources/styles/indexstyle.ist, intoc]
\indexsetup{level=\chapter*,toclevel=chapter}

%Start des Inhalts
\begin{document}

%Notwendiger Workaround
\pagenumbering{alph}

%Deckblatt erzeugen
\maketitle

\pagenumbering{Roman}

\section{Schaltverhalten der Undeland-Beschaltung}


\subsection{Schaltverhalten des Schalters T1}
\begin{figure}
\centering
\includegraphics[width=0.8\textwidth]{figures/s1_v_i.eps}
\caption{Spannung und Stromverlauf des Schalters T1}
\label{fig:s1_v_i}
\end{figure}

\subsection{Spannung an den WS-Klemmen}

\begin{figure}
\centering
\includegraphics[width=0.8\textwidth]{figures/ws_v.eps}
\caption{Spannung an den WS-Klemmen}
\label{fig:ws_v}
\end{figure}

\subsection{Berechnung der differenziellen Strom und Spannung an T1}

Über den Spannungsabfall der Induktivität $L_\mathrm{E}=\qty[]{12}{\micro H}$ und die GS-Spannung $U_\mathrm{D}=\qty[]{2700}{V}$  lässt sich das $\frac{di}{dt}$
\begin{align}
    U_\mathrm{L}&=L\frac{di}{dt}\\
    \frac{di}{dt}&=\frac{U_\mathrm{L}}{L}\\ 
    \frac{di}{dt}&=\frac{\qty[]{2700}{V}}{\qty[]{12}{ \micro H}}\\ 
    \frac{di}{dt}&=\qty[]{225}{A/\micro s}\\
\end{align}
bestimmen.

Über den Laststrom $\hat I=\qty[]{1.3}{kA}$ und die Kapazität $C_s=\qty[]{3}{\micro F}$ lässt sich das $\frac{du}{dt}$

\begin{align}
    I_\mathrm{C}&=C\frac{du}{dt}\\
    \frac{du}{dt}&=\frac{I_\mathrm{C}}{C}\\ 
    \frac{du}{dt}&=\frac{\qty[]{1.3}{kA}}{\qty[]{3}{ \micro F}}\\ 
    \frac{du}{dt}&=\qty[]{433}{V/\micro s}\\
\end{align}

\subsection{Bestimmung der differenziellen Strom und Spannung an T1 mit Hilfe der Simualtion}

\begin{figure}
\centering
\includegraphics[width=0.8\textwidth]{figures/dt.eps}
\caption{Differenzielle Spannung und Strom an T1 Simulation}
\label{fig:dt}
\end{figure}

Aus der Simulation werden über Marker die von Spannung und Strom in Abhängigkeit der Zeit, für das Ein- und Ausschalten abgelesen.

\begin{equation}
    \frac{du}{dt}=\frac{u_2-u_1}{t_2-t_1}=\frac{\qty[]{2938}{V}-\qty[]{3.146}{V}}{\qty[]{15.11967}{ms}-\qty[]{15.1124}{ms}}=\qty[]{403.7}{V/\micro s}
\end{equation}

\begin{equation}
    \frac{di}{dt}=\frac{i_2-i_1}{t_2-t_1}=\frac{\qty[]{2538}{A}-\qty[]{0.0270}{A}}{\qty[]{11.2554}{ms}-\qty[]{11.241}{ms}}=\qty[]{176.25}{A/\micro s}
\end{equation}

Bei der Berechnung der Werte wurde mit der maximalen Spannung sowie dem maximalen Strom gerechnet. Da in der Simulation nicht im Maximalwert geschaltet wird, erklärt sich die geringe Abweichung.

Die maximale Spannung am Schalter ist hierbei $U_\mathrm{max}=\qty[]{2938}{V}$.

\subsection{Effektivwert der Ströme in $C_\mathrm{s},\,D_\mathrm{S1},\,D_\mathrm{S2}$ und $R_\mathrm{Cl}$}

\begin{align}
    I_\mathrm{Cs RMS}&=\qty[]{58.716}{A}\\ 
    I_\mathrm{DS1 RMS}&=\qty[]{159.89}{A}\\
    I_\mathrm{DS2 RMS}&=\qty[]{158.47}{A}\\ 
    I_\mathrm{RCl RMS}&=\qty[]{62.139}{A}\\
\end{align}

\begin{figure}[h]
\centering
\includegraphics[width=0.8\textwidth]{figures/ucl.eps}
\caption{Spannung an $R_\mathrm{Cl}$}
\label{fig:ucl}
\end{figure} \clearpage
%\include{chapter/Abstract} \clearpage

%Inhaltsverzeichnis
\tableofcontents \newpage

%Hauptteil
\pagenumbering{arabic}

%\input{chapter/Einleitung} \clearpage
%\input{chapter/} \clearpage

%Anhang
\pagenumbering{Alph}

%Abbildungsverzeichnis
\listoffigures \clearpage
%Tabellenverzeichnis
%\listoftables \clearpage
%Quelltextverzeichnis
%\lstlistoflistings \clearpage
%Stichwortverzeichnis
%\printindex \clearpage
%Glossar
%\printglossary[title={Glossar}] \clearpage
%Abkürzungsverzeichnis
%\printglossary[style=dottedlocations,type=\acronymtype,title={Abkürzungsverzeichnis}] \clearpage

%Literaturverzeichnisse (getrennt nach Stichwort)
\printbibliography[heading=bibintoc, keyword={book}, title={Literaturverzeichnis}]\clearpage
\printbibliography[heading=bibintoc, keyword={online}, title={Onlinequellen}]\clearpage
\printbibliography[heading=bibintoc, keyword={image}, title={Bildquellen}]\clearpage

% Anhang
%\input{chapter/Anhang}

% Eigenständigkeitserklärung
%\input{chapter/Eigenstaendigkeitserklaerung}

\end{document}